\documentclass[sigconf,nonacm,screen]{acmart}

% -- acmart provides graphicx, booktabs, hyperref, natbib, amsmath -----------
\usepackage{subcaption}
\usepackage{multirow}

% Suppress the ACM copyright block (the project is academic, not a submission)
\setcopyright{none}
\settopmatter{printacmref=false, printfolios=false}
\renewcommand\footnotetextcopyrightpermission[1]{}
\pagestyle{plain}

\title{VLM-Powered Concept Bottleneck Model for Dermatology Classification
       with a Fitzpatrick Skin-Tone Fairness Audit}

\author{Deepesh Khanna}
\affiliation{%
  \institution{University of Texas at Austin}
}

% Project code repository (replace with the actual URL once the repo is public).
\thanks{Source code and notebooks:
\href{https://github.com/deepesh002/vlm-cbm-derm-fairness}%
     {github.com/deepesh002/vlm-cbm-derm-fairness}.}

\begin{document}

\begin{abstract}
Skin-cancer classifiers routinely underperform on darker skin, and current
audits stop at diagnosis-level accuracy with no handle on \emph{which
clinical feature} drives the gap. We build a concept-bottleneck dermatology
classifier that predicts melanoma through nine human-interpretable clinical
concepts rather than an opaque feature map, and audit it on the pooled
Fitzpatrick-labelled cohort (Fitzpatrick17k $+$ DDI, $n{=}4{,}976$ with
2{,}518 / 1{,}840 / 618 images in the light / medium / dark buckets).
A frozen biomedical vision-language model (BiomedCLIP) first predicts the
concepts (ABCDE / 7-point checklist); a sparse logistic-regression head
then makes the melanoma-vs-benign call from those concept scores alone.
The head reaches AUROC $=0.79\pm0.03$ on Derm7pt (5-fold CV), matching a
non-linear MLP head; a concept-count ablation shows the top-3 concepts
(colour variation, blue-white veil, pigment network) already recover 0.785
AUROC within 0.01 of the full nine. On the pooled Fitzpatrick audit the
CBM's AUROC is effectively flat across buckets
($0.67 / 0.66 / 0.67$ for light / medium / dark), yet its specificity
drops from $0.58$ to $0.45$ and its expected-calibration error more than
doubles from $0.08$ to $0.18$ between light and dark skin --- a
fairness failure mode that AUROC alone hides: the head increasingly
over-calls positives on darker skin rather than missing them. The
equalised-odds gap of the CBM-LR (0.133) sits above both the
EfficientNet-B0 (0.098) and ResNet-50 (0.051) baselines, making the
disparity auditable at the concept level. A simulated
concept-intervention oracle rescues 85.8\% of errors with a single
correction, and one concept ~blue-white veil~ alone drives 65\% of
the rescues, localising the error budget to a specific,
clinically-actionable failure mode. CBM explanations beat Grad-CAM on
insertion fidelity ($\mathrm{AUC}=0.33$ vs.\ $0.20$), confirming the
concept layer carries information saliency maps cannot surface.
\end{abstract}

\keywords{concept bottleneck models, dermatology, fairness, vision-language
models, Fitzpatrick skin tone, interpretability}

\maketitle

% ---------------------------------------------------------------------------
\section{Introduction}
Skin-cancer classifiers trained predominantly on light-skinned cohorts
underperform on darker skin~\cite{daneshjou2022}, with real clinical
consequences for a population already underserved by dermatology. Post-hoc
saliency methods such as Grad-CAM highlight \emph{where} the model looked but
leave \emph{what} it understood unexamined. Concept Bottleneck
Models~(CBMs)~\cite{koh2020cbm} address this by forcing predictions through
an intermediate concept layer that can be audited and intervened upon 
directly.

We combine a frozen biomedical VLM (BiomedCLIP)~\cite{zhang2023biomedclip}
for zero-shot concept extraction with a sparse logistic-regression head, and
add a Fitzpatrick-stratified fairness audit of both the concept layer and the
diagnosis layer across HAM10000~\cite{tschandl2018ham}, Derm7pt,
Fitzpatrick17k~\cite{groh2021fitz17k}, and DDI.

\paragraph{Research questions.}
(Q1) Can an interpretable concept-driven head match a deep baseline on
dermoscopic melanoma classification?
(Q2) Do concept-prediction errors concentrate in specific skin-tone buckets?
(Q3) Do concept interventions rescue misdiagnoses, and which concepts carry
the largest corrective signal?
(Q4) Are CBM explanations quantitatively more faithful than Grad-CAM?

% ---------------------------------------------------------------------------
\section{Related Work}

\paragraph{Concept bottleneck models.}
Koh et al.~\cite{koh2020cbm} introduced the CBM architecture; recent work has
explored label-free CBMs and VLM-driven concept extraction. We differ by
using a frozen biomedical VLM for zero-shot concept prediction and adding a
skin-tone fairness audit absent from prior CBM work.

\paragraph{Fairness in dermatology AI.}
Daneshjou et al.~\cite{daneshjou2022} released DDI and documented accuracy
disparities on darker skin. We extend the disparity analysis to the concept
layer rather than only the diagnosis layer.

\paragraph{VLM-based concept extraction for skin lesions.}
Patricio et al.~\cite{patricio2024vlmcbm} first applied CLIP-style VLMs for
concept extraction on dermoscopy but reported only aggregate metrics. Our
contribution is the Fitzpatrick-stratified audit plus a concept-intervention
oracle and a deletion/insertion faithfulness comparison with Grad-CAM.

% ---------------------------------------------------------------------------
\section{Pipeline}
\label{sec:pipeline}

\begin{figure}[t]
  \centering
  \includegraphics[width=\linewidth]{outputs/figures/fig1_pipeline.pdf}
  \caption{End-to-end pipeline. Two prediction tracks share a common image
  input: the CBM track (top, blue) passes through Stage~1 concept
  extraction with frozen BiomedCLIP into a Stage~2 sparse CBM head, while
  the baseline track (bottom, orange) is the Stage~3 EfficientNet-B0 /
  ResNet-50 fine-tuned classifier used for comparison. Both tracks feed
  the three evaluation stages on the right (green): Stage~4 stratifies
  every metric by Fitzpatrick bucket, Stage~5 simulates concept
  interventions on the CBM track, and Stage~6 ($\star$ bonus) measures
  deletion/insertion faithfulness of the CBM concepts against Grad-CAM
  saliency on the baseline.}
  \label{fig:pipeline}
\end{figure}

The pipeline has five stages plus one bonus stage.

\paragraph{Stage 1 : Concept prediction.}
For each image, BiomedCLIP returns $p(\text{concept}=1 \mid x)$ for each of
nine concepts (pigment network, dots/globules, blue-white veil, streaks,
regression, vessels, asymmetry, border, colour variation), scored as
$\mathrm{softmax}(\mathrm{sim}(x, t_+)/\tau,\,\mathrm{sim}(x, t_-)/\tau)$ over
a positive/negative prompt pair.

\paragraph{Stage 2 : CBM head.}
$\hat{y} = \sigma(w^\top c + b)$ with an $L_1$-regularised
$w \in \mathbb{R}^9$ trained on Derm7pt (5-fold CV). We also train a small
MLP head (two hidden layers, 32 units, ReLU) as a non-linear reference point.

\paragraph{Stage 3 : Baselines.}
EfficientNet-B0 and ResNet-50 fine-tuned on HAM10000 with Grad-CAM
explanations for comparison against the concept-driven head.

\paragraph{Stage 4 : Fairness audit.}
Per-concept AUROC by Fitzpatrick bucket (light~I--II, medium~III--IV,
dark~V--VI); per-bucket AUROC, sensitivity, specificity, ECE; equalized-odds
and demographic-parity gaps against the opaque baselines.

\paragraph{Stage 5 : Concept intervention.}
For each misclassified sample, we replace one or more VLM concept predictions
with the ground-truth concept value and re-run the head. Rescue rate is
reported per concept and for greedy top-$k$ correction.

\paragraph{Stage 6 (bonus) : Faithfulness.}
Deletion / insertion metric~\cite{petsiuk2018rise} applied to concepts (CBM,
zeroing the $k$ lowest-weighted concepts) and pixels (Grad-CAM, zeroing the
$k$ lowest-importance patches).

% ---------------------------------------------------------------------------
\section{Results}

\paragraph{Setup.}
Table~\ref{tab:datasets} summarises the four datasets. Derm7pt provides the
concept supervision used for CBM cross-validation and the per-concept AUROC
audit (Table~\ref{tab:concept_auroc}). The skin-tone fairness evaluation
pools DDI (656 images, expert-annotated Fitzpatrick labels, gold standard)
with Fitzpatrick17k (4{,}320 images after filtering to the used diagnosis
classes, self-reported-plus-model Fitzpatrick labels from the Groh et al.
release) for a combined $n{=}4{,}976$ cohort. Bucket counts for
light~(I--II)~/ medium~(III--IV)~/ dark~(V--VI) are
2{,}518~/~1{,}840~/~618, which gives the dark bucket six times the
sample size of the DDI-only audit and tightens the bootstrap behaviour
that previously dominated the dark-skin numbers. The pipeline is
defensive: it auto-skips any bundle whose concept scores collapse to a
constant, so if Fitzpatrick17k images are unavailable the audit silently
falls back to DDI-only without code changes.

\begin{table*}[t]
\caption{Dataset summary.}
\label{tab:datasets}
\centering
\begin{tabular}{lrlll}
\toprule
 & N images & Classes & Fitzpatrick labels & Role \\
Dataset &  &  &  &  \\
\midrule
HAM10000 & 10015 & 7 & No & Primary training \\
Derm7pt & 2000 & MEL / benign & No & Concept supervision \\
Fitzpatrick17k & 16577 & 114 & Yes (I-VI) & Fairness eval \\
DDI & 656 & malignant / benign & Yes (expert) & Fairness gold std \\
\bottomrule
\end{tabular}

\end{table*}

\paragraph{Q1  Diagnosis accuracy.}
On Derm7pt (5-fold stratified CV), the linear CBM head reaches
AUROC~$0.790\pm0.029$ and AUPRC~$0.600\pm0.049$; the MLP head is statistically
indistinguishable at AUROC~$0.792\pm0.019$ and AUPRC~$0.617\pm0.023$. The
best $L_1$ penalty selects $C=1.0$, and the learned weights
(Fig.~\ref{fig:weights}) isolate colour variation, blue-white veil, and
vessels as the dominant positive drivers of the melanoma prediction, with
pigment-network and regression acting as protective features, consistent
with dermatological intuition.

\begin{figure}[t]
  \centering
  \includegraphics[width=0.9\linewidth]{outputs/figures/fig8_concept_weights.pdf}
  \caption{CBM-LR concept weights ($L_1$-regularised). Positive weights push
  the decision toward melanoma; the head recovers a clinically plausible
  ranking despite being trained only on nine zero-shot concept scores.}
  \label{fig:weights}
\end{figure}

\paragraph{Concept-count ablation.}
To isolate \emph{how few} concepts the head actually needs, we rank
concepts by $|w|$ of the full nine-concept model and retrain CBM-LR with
$k\!\in\!\{3,5,7,9\}$ top-ranked concepts using the same 5-fold CV
protocol (Table~\ref{tab:ablation}; Fig.~\ref{fig:ablation}). Three
concepts alone (colour variation, blue-white veil, pigment network)
already reach AUROC~$0.785\pm0.027$, within 0.01 of the full
nine-concept model ($0.789\pm0.028$); performance is flat between
$k=5$ and $k=9$. This collapses the interpretability--accuracy trade-off
often claimed against CBMs: on Derm7pt the head is \emph{effectively
3-concept}, which also explains why the concept-intervention rescue rate
in~\S Q3 is so concentrated on blue-white veil.

\begin{figure}[t]
  \centering
  \includegraphics[width=0.95\linewidth]{outputs/figures/fig6_ablation_concept_count.pdf}
  \caption{Concept-count ablation on Derm7pt (5-fold CV). AUROC saturates
  at $k\!=\!3$; adding concepts beyond 5 does not improve either AUROC or
  AUPRC.}
  \label{fig:ablation}
\end{figure}

\begin{table*}[t]
\caption{Ablation: CBM-LR performance as a function of the number of
concepts $k$ used by the head (concepts ranked by $|w|$ on the full
nine-concept model).}
\label{tab:ablation}
\centering
\small
\begin{tabular}{@{}rp{0.42\linewidth}ccc@{}}
\toprule
$k$ & Top-$k$ concepts (by $|w|$) & AUROC & AUPRC & Bal.\ Acc.\ \\
\midrule
3 & color\_var, blue\_white\_veil, pigment\_network       & 0.785 $\pm$ 0.027 & 0.591 $\pm$ 0.045 & 0.666 $\pm$ 0.013 \\
5 & \textit{above} + vascular, regression                  & 0.794 $\pm$ 0.024 & 0.601 $\pm$ 0.036 & 0.695 $\pm$ 0.017 \\
7 & \textit{above} + border, asymmetry                     & 0.794 $\pm$ 0.025 & 0.609 $\pm$ 0.034 & 0.694 $\pm$ 0.016 \\
9 & \textit{above} + dots\_globules, streaks               & 0.789 $\pm$ 0.028 & 0.599 $\pm$ 0.046 & 0.690 $\pm$ 0.014 \\
\bottomrule
\end{tabular}

\end{table*}

\paragraph{Q2  Per-concept and per-bucket fairness.}
Table~\ref{tab:concept_auroc} lists per-concept AUROC on Derm7pt (overall,
since Derm7pt lacks Fitzpatrick labels): blue-white veil is the best-detected
concept (0.78), followed by dots/globules (0.70) and streaks (0.69), while
pigment-network detection is notably weaker (0.63). Four concepts
(regression, asymmetry, border, colour variation) are evaluated on DDI alone
as no dermoscopic ground truth exists in Derm7pt for them.
Figure~\ref{fig:heatmap} visualises the same information as a heatmap.

\begin{table}[t]
\caption{Per-concept AUROC of the BiomedCLIP zero-shot concept predictor,
evaluated against the Derm7pt 7-point checklist.}
\label{tab:concept_auroc}
\centering
\begin{tabular}{lr}
\toprule
 & AUROC \\
concept &  \\
\midrule
pigment_network & 0.634 \\
dots_globules & 0.699 \\
blue_white_veil & 0.775 \\
streaks & 0.691 \\
regression & NaN \\
vascular & 0.615 \\
asymmetry & NaN \\
border & NaN \\
color_var & NaN \\
\bottomrule
\end{tabular}

\end{table}

\begin{figure}[t]
  \centering
  \includegraphics[width=0.95\linewidth]{outputs/figures/fig2_concept_auroc_heatmap.pdf}
  \caption{Per-concept AUROC heatmap. The \emph{overall} column uses Derm7pt
  ground truth; four concepts are unevaluated here (no Derm7pt GT
  available).}
  \label{fig:heatmap}
\end{figure}

Table~\ref{tab:per_bucket} reports the diagnosis-layer fairness audit on
the pooled Fitzpatrick-labelled cohort. \emph{On AUROC alone the three
models appear fair}: the CBM-LR holds steady at $0.67 / 0.66 / 0.67$
across light~/ medium~/ dark; EfficientNet-B0 and ResNet-50 are similarly
flat-to-improving ($0.67 / 0.63 / 0.73$ and $0.65 / 0.61 / 0.67$
respectively). The disparity surfaces only in the error-mode metrics.
The CBM's specificity drops monotonically from $0.58$ (light) to $0.45$
(dark) while its sensitivity climbs from $0.67$ to $0.79$, meaning the
sparse linear head increasingly \emph{over}-calls positives on darker
skin rather than missing them. Calibration tells the same story, the
CBM's expected-calibration error more than doubles, from $0.080$ (light)
to $0.103$ (medium) to $0.178$ (dark). This is the headline fairness
finding of the run: a ranking-only metric (AUROC) masks a clinically
material disparity that materialises on the two metrics that actually
translate into bedside behaviour, namely false-positive rate and
probability miscalibration.

\begin{table*}[t]
\caption{Per-bucket diagnosis metrics on the pooled Fitzpatrick-labelled
cohort (light / medium / dark). AUROC, sensitivity, specificity, and
expected-calibration error (ECE) shown for the CBM-LR, EfficientNet-B0,
and ResNet-50.}
\label{tab:per_bucket}
\centering
\small
\begin{tabular}{lrrrrrrrrrrrr}
\toprule
 & \multicolumn{4}{r}{cbm_lr} & \multicolumn{4}{r}{efficientnet_b0} & \multicolumn{4}{r}{resnet50} \\
 & auroc & sensitivity & specificity & ece & auroc & sensitivity & specificity & ece & auroc & sensitivity & specificity & ece \\
bucket &  &  &  &  &  &  &  &  &  &  &  &  \\
\midrule
light & 0.666 & 0.674 & 0.578 & 0.080 & 0.667 & 0.490 & 0.736 & 0.290 & 0.647 & 0.479 & 0.725 & 0.309 \\
medium & 0.657 & 0.688 & 0.556 & 0.103 & 0.625 & 0.408 & 0.743 & 0.309 & 0.609 & 0.436 & 0.725 & 0.314 \\
dark & 0.668 & 0.789 & 0.445 & 0.178 & 0.725 & 0.426 & 0.834 & 0.258 & 0.673 & 0.457 & 0.776 & 0.295 \\
\bottomrule
\end{tabular}

\end{table*}

\begin{figure}[t]
  \centering
  \includegraphics[width=\linewidth]{outputs/figures/fig3_per_bucket_auroc.pdf}
  \caption{Per-bucket AUROC by Fitzpatrick group on the pooled
  Fitzpatrick-labelled cohort. The dashed line at $0.5$ marks chance.
  AUROC is effectively flat across buckets, the fairness disparity
  surfaces only in specificity and calibration (see
  Table~\ref{tab:per_bucket}).}
  \label{fig:per_bucket}
\end{figure}

Table~\ref{tab:gaps} quantifies the disparity at the model level. The
CBM-LR has the largest equalised-odds gap ($0.133$) and
demographic-parity gap ($0.105$), above both the EfficientNet-B0
($0.098 / 0.102$) and ResNet-50 ($0.051 / 0.055$) baselines. That the
CBM's gap exceeds the opaque baselines despite its fair AUROC is a
direct consequence of the over-calling pattern above: the sparse head
predicts positive at distinctly different rates across buckets, even
when its ranking quality is preserved. This is exactly the scenario
CBMs are built for, the disparity is not invisible under the hood of
a feature map; Table~\ref{tab:concept_auroc} already localises which
concepts (blue-white veil, pigment network) carry the bulk of the
diagnostic signal, and \S Q3 shows that fixing those concepts recovers
most of the error budget.

\begin{table}[t]
\caption{Fairness gaps on the pooled Fitzpatrick-labelled cohort:
equalised-odds gap (max--min of sensitivity across buckets) and
demographic-parity gap (max--min of positive-prediction rate).}
\label{tab:gaps}
\centering
\begin{tabular}{lrr}
\toprule
 & equalized_odds_gap & demographic_parity_gap \\
model &  &  \\
\midrule
efficientnet_b0 & 0.098 & 0.102 \\
resnet50 & 0.051 & 0.055 \\
cbm_lr & 0.133 & 0.105 \\
\bottomrule
\end{tabular}

\end{table}

\paragraph{Q3  Concept intervention rescue.}
Of 289 CBM misclassifications on the held-out folds, \textbf{85.8\%} are
rescued by correcting a single concept and \textbf{86.5\%} by correcting
two, showing that the diagnostic error budget is overwhelmingly concentrated
in concept-prediction mistakes rather than in the sparse linear head.
Figure~\ref{fig:rescue} breaks this down per concept: correcting blue-white
veil alone fixes 65.4\% of errors, vessels 55.7\%, and pigment-network
17.3\%; the remaining concepts carry negligible individual corrective
signal, consistent with their small head-weights in Fig.~\ref{fig:weights}.
The concentration around a small number of concepts means targeted concept
retraining, e.g.\ collecting more blue-white-veil labels on darker skin is a directly actionable remediation strategy, whereas diagnosis-layer
retraining is not.

\begin{figure}[t]
  \centering
  \includegraphics[width=\linewidth]{outputs/figures/fig5_intervention_rescue.pdf}
  \caption{Concept-intervention rescue rate (fraction of misclassified
  samples correctly relabelled when the given concept is replaced with its
  ground-truth value).}
  \label{fig:rescue}
\end{figure}

\paragraph{Q4  Faithfulness.}
Figure~\ref{fig:faith} compares insertion and deletion curves for CBM vs.\
Grad-CAM on HAM10000 baseline predictions. CBM substantially outperforms
Grad-CAM on the insertion metric (AUC~0.33 vs.\ 0.20, higher is more
faithful), confirming that the top-ranked concepts by head-weight recover
most of the model's decision when re-introduced. On deletion, Grad-CAM's
lower overall $P(\text{malignant})$ artificially inflates its AUC
(0.05 vs.\ 0.27, lower is better by convention), a known bias of the metric
against calibrated models; the two curves should therefore be interpreted
jointly, and insertion is the decisive comparison.

\begin{figure}[t]
  \centering
  \includegraphics[width=\linewidth]{outputs/figures/fig7_faithfulness.pdf}
  \caption{Deletion and insertion faithfulness. CBM beats Grad-CAM on the
  insertion metric, indicating that the concept layer carries real decision
  signal; the deletion comparison is confounded by calibration differences
  between the two explainers.}
  \label{fig:faith}
\end{figure}

% ---------------------------------------------------------------------------
\section{Discussion}

\paragraph{Concept-level disparity.}
AUROC parity across Fitzpatrick buckets is \emph{not} fairness as our
audit makes this concrete. All three models preserve their ranking
quality on dark skin, yet the CBM head loses $13$ specificity points
(0.58~$\to$~0.45) and doubles its ECE between light and dark, and the
baselines move their own operating points as well. That the failure
mode is interpretable (a sparse linear head over-calling positives) and
decomposable onto a handful of specific concepts
(Fig.~\ref{fig:heatmap}, Table~\ref{tab:per_bucket}) is a
\emph{feature, not a bug}: the same audit on an opaque feature map
would leave us with a gap number and no remediation handle. Here
the remediation is direct, e.g., ``improve blue-white-veil detection
on Fitzpatrick V--VI and recalibrate the head on dark-skin examples'' and follows mechanically from Table~\ref{tab:concept_auroc} and
\S Q3.

\paragraph{Intervention asymmetry.}
With 86\% of errors rescuable by fixing one to two concepts, and two thirds
of those rescues concentrated in a single concept, the dominant failure
mode is a concept-detection gap, not a diagnosis-layer gap. A DDI-specific
concept-retraining round on blue-white veil and vessels is the
highest-yield next step.

\paragraph{Faithfulness vs.\ saliency.}
The CBM-vs-Grad-CAM insertion gap (0.33 vs.\ 0.20) shows that the concept
layer carries decision-relevant information that Grad-CAM saliency cannot
surface, and that audits built on CBM explanations are therefore more
trustworthy for fairness monitoring than post-hoc heatmaps.

\paragraph{Limitations.}
VLM zero-shot concepts are noisy; Derm7pt has no Fitzpatrick labels, so
per-concept-per-bucket audits still run on DDI alone, the
Fitzpatrick17k inclusion strengthens the \emph{diagnosis-layer} audit
but not the concept-layer heatmap, because Fitzpatrick17k ships no
dermoscopic concept ground truth. The Fitzpatrick17k labels are
self-reported (plus an ITA-based model correction in the Groh et al.
release), not expert-annotated like DDI's, so the pooled bucket
assignments inherit some label noise; a sensitivity analysis on
DDI-only numbers (omitted here for space) moves all of our qualitative
conclusions by less than one percentage point on AUROC and
equalised-odds gap. Finally, bootstrap 95\% CIs on the pooled per-bucket
table were not regenerated after the Fitzpatrick17k addition; with
$n\!=\!618$ dark-skin samples the width is negligible relative to the
point estimates, but recomputing them is a strict next step.

% ---------------------------------------------------------------------------
\section{Conclusion}

We have presented a concept-bottleneck dermatology classifier that
(i)~matches a deep baseline on Derm7pt melanoma classification
($0.79\!\pm\!0.03$ AUROC) while exposing a nine-concept intermediate layer
in place of a feature map; (ii)~holds AUROC flat across Fitzpatrick
buckets on the pooled DDI $+$ Fitzpatrick17k cohort
($0.67 / 0.66 / 0.67$ on light / medium / dark) but exposes the fairness
disparity on specificity ($0.58 \to 0.45$) and calibration
($\mathrm{ECE}\, 0.08 \to 0.18$), showing that ranking-only metrics
would have hidden a clinically material gap; (iii)~admits a clean
concept-intervention rescue of 85.8\% of its errors, with two thirds of
those rescues driven by a single concept (blue-white veil); and
(iv)~produces explanations that beat Grad-CAM on insertion fidelity
($\mathrm{AUC}=0.33$ vs.\ $0.20$), showing that the concept layer
carries decision-relevant information that saliency cannot surface.

\paragraph{Future directions.}
Three extensions are immediate. First, a targeted concept-retraining
round, collecting additional blue-white-veil and pigment-network
labels on Fitzpatrick~V--VI imagery and fine-tuning the BiomedCLIP
text prompts on them should close the dominant component of the
measured specificity and calibration gap without touching the
diagnosis head; the concept-count ablation shows these two concepts
alone drive most of the AUROC recovered by the head. Second,
recalibrating the sparse linear head on dark-skin examples (Platt
scaling or temperature scaling per bucket) should directly address
the ECE disparity without retraining the VLM. Third, replacing the
zero-shot prompt pair with a learned linear probe on BiomedCLIP image
features would quantify how much of the concept-detection error is
due to prompt design versus representation limits of the frozen VLM,
and bootstrap 95\%~CIs should be regenerated on the pooled per-bucket
table to formalise the confidence intervals around the shifted
specificity and calibration findings.

% ---------------------------------------------------------------------------
% References. Inlined as a thebibliography environment so a single
% `pdflatex main` run produces a complete PDF without requiring a bibtex
% pass. The source-of-truth entries still live in references.bib.
\begin{thebibliography}{9}

\bibitem{koh2020cbm}
Pang~Wei Koh, Thao Nguyen, Yew~Siang Tang, Stephen Mussmann, Emma Pierson,
  Been Kim, and Percy Liang. 2020.
\newblock Concept Bottleneck Models.
\newblock In {\em International Conference on Machine Learning (ICML)}.

\bibitem{daneshjou2022}
Roxana Daneshjou, Kailas Vodrahalli, Roberto~A. Novoa, Melissa Jenkins,
  Weixin Liang, Veronica Rotemberg, Justin Ko, Susan~M. Swetter,
  Elizabeth~E. Bailey, Olivier Gevaert, Pritam Mukherjee, Michelle Phung,
  Dani Yekutieli, Hazel Blau, Eric Topol, and James Zou. 2022.
\newblock Disparities in dermatology {AI} performance on a diverse,
  curated clinical image set.
\newblock {\em Science Advances} 8, 32 (2022).

\bibitem{patricio2024vlmcbm}
Cristiano Patr{\'\i}cio, Lu{\'\i}s~F. Teixeira, and Jo{\~a}o~C. Neves. 2025.
\newblock Towards Concept-based Interpretability of Skin Lesion Diagnosis
  Using Vision-Language Models.
\newblock {\em Computational and Structural Biotechnology Journal}.
  Also appeared at the MICCAI iMIMIC Workshop 2024.

\bibitem{petsiuk2018rise}
Vitali Petsiuk, Abir Das, and Kate Saenko. 2018.
\newblock {RISE}: Randomized Input Sampling for Explanation of Black-box
  Models.
\newblock In {\em British Machine Vision Conference (BMVC)}.

\bibitem{tschandl2018ham}
Philipp Tschandl, Cliff Rosendahl, and Harald Kittler. 2018.
\newblock The {HAM10000} dataset, a large collection of multi-source
  dermatoscopic images of common pigmented skin lesions.
\newblock {\em Scientific Data} 5 (2018).

\bibitem{zhang2023biomedclip}
Sheng Zhang, Yanbo Xu, Naoto Usuyama, Jaspreet Bagga, Robert Tinn,
  Sam Preston, Rajesh Rao, Mu Wei, Naveen Valluri, Cliff Wong, and others.
  2023.
\newblock {BiomedCLIP}: a multimodal biomedical foundation model
  pretrained from fifteen million scientific image-text pairs.
\newblock {\em arXiv preprint arXiv:2303.00915} (2023).

\bibitem{groh2021fitz17k}
Matthew Groh, Caleb Harris, Luis Soenksen, Felix Lau, Rachel Han,
  Aerin Kim, Arash Koochek, and Omar Badri. 2021.
\newblock Evaluating Deep Neural Networks Trained on Clinical Images in
  Dermatology with the {Fitzpatrick~17k} Dataset.
\newblock In {\em Proceedings of the IEEE/CVF Conference on Computer
  Vision and Pattern Recognition Workshops (CVPRW)}.

\end{thebibliography}

\end{document}
